\section{风险、边界与后续计划}

\subsection{当前必须诚实说明的边界}

一个比赛项目要具备说服力，关键不只是展示完成了什么，还要明确说明尚未完成什么。
当前项目最重要的边界包括：

\begin{itemize}
    \item 冻结比赛主线 ISA 仍是 \texttt{RV32I + Zicsr}；
    \item 扩展 \texttt{full-ui} 使用了 \texttt{zifencei} 覆盖矩阵，但不应被误写成
    ``冻结比赛 ISA 已升级''；
    \item FPGA 当前是 pre-board closure，不是实板 signoff；
    \item 仍未完成的板级工作都属于外部阻塞，而不是本机能力内尚未做的基本任务。
\end{itemize}

\subsection{当前主要风险}

\begin{table}[H]
    \centering
    \caption{当前风险与控制策略}
    \begin{tabular}{p{0.25\textwidth} p{0.28\textwidth} p{0.31\textwidth}}
        \toprule
        风险点 & 风险描述 & 当前控制策略 \\
        \midrule
        文档口径漂移 & 结果、脚本和汇报材料引用值不一致 & 统一由冻结基线与证据索引维护正式口径 \\
        高风险优化误入主线 & 局部提分但打坏整体正确性 & 单变量实验、先过 guardrail、再跑完整回归 \\
        ISA 边界误读 & 将扩展 UI 覆盖误认为主线 ISA 升级 & 在技术文档和本说明书中明确写清边界 \\
        FPGA 状态误报 & 将 pre-board 写成已上板闭环 & 在检查清单和交付文档中保留外部阻塞说明 \\
        \bottomrule
    \end{tabular}
\end{table}

\subsection{近期优化结论}

项目并没有停留在 ``跑出一版结果就停止'' 的状态，而是持续做了多轮性能探索。
但这些探索并非都值得保留。当前已经明确的结论包括：

\begin{enumerate}
    \item request-side 与 queue/reuse 微调方向已基本被证据排除；
    \item \texttt{FQ-06A} quick-screen 虽然诊断全绿，但 CoreMark short 零收益，因此未保留；
    \item 2026-04-09 的 split profile 说明 redirect 成本主要来自 taken branch，尤其是
    \texttt{bne} 与 \texttt{beq}；
    \item 一个 ``decode-stage early JAL redirect'' 试验已经被 guardrail 否定，
    因为它打破了当前隐含的控制气泡平衡，风险高于收益。
\end{enumerate}

这意味着项目后续如果继续优化，不应该回到已经被证伪的 request/queue 微调，
而应从更严谨的控制流成本分解出发，重新建立新的单变量假设。

\subsection{后续计划}

从比赛项目管理视角看，后续工作可以划分为三个方向：

\begin{enumerate}
    \item \textbf{持续维护交付材料}：
    保持说明书、技术文档、交接说明和证据索引一致，确保任何时点都能直接交接。
    \item \textbf{继续做证据驱动的优化}：
    若后续继续做性能实验，应优先围绕 taken branch redirect 成本展开，
    并保持 ``单变量、小步、可回归'' 原则。
    \item \textbf{等待外部条件后完成实板闭环}：
    当实体板卡与最终约束条件具备后，再补做 UART/LED/boot banner 与板级 signoff。
\end{enumerate}

\subsection{本说明书在后续阶段的作用}

当前 LaTeX 说明书不仅是初赛提交文档，也是后续比赛阶段的持续维护载体：

\begin{itemize}
    \item 若新增架构图、波形图或 Vivado 截图，可直接插入本说明书；
    \item 若未来加入新的保留优化，可在不打乱结构的前提下更新对应章节；
    \item 若赛方要求更正式的 PDF 包、目录或附录结构，可以直接基于本版本扩写。
\end{itemize}

因此，本说明书的目标不是 ``一次性写完''，而是把当前工程真实状态转化为可持续维护的比赛文档基线。

\clearpage
